\documentclass[11pt,letterpaper]{article}

%%% LANGUAGE: Spanish
\usepackage[spanish,english]{babel}
\addto\shorthandsspanish{\spanishdeactivate{~<>}}

%%% ENCODING
\usepackage[utf8]{inputenc}
\usepackage{microtype}
\usepackage[T1]{fontenc}
\usepackage{lmodern}
\usepackage{amsmath}
\usepackage{textcomp}

\usepackage{graphicx}
\usepackage{epstopdf}
\usepackage{epsfig}

%%
%%%%GENERAL INFROMATION%%%%
%%

\newcommand{\theCampus}{Campus Vitacura}
\newcommand{\theAssignment}{SIG}
\newcommand{\therevision}{DV}

%%% PACKAGES
\usepackage{booktabs} 	% for much better looking tables
\usepackage{array} 		% for better arrays (eg matrices) in maths
\usepackage{paralist} 	% very flexible & customizable lists (eg. enumerate/itemize, etc.)
\usepackage{verbatim} 	% adds environment for commenting out blocks of text & for better verbatim
\usepackage{subfigure} % make it possible to include more than one captioned figure/table in a single float
\usepackage{xcolor}
\usepackage{listings}  % bloques de código SQL

% FECHA
\usepackage{datetime}
\newdateformat{monthyeardate}{%
  \monthname[\THEMONTH], \THEYEAR}

%%% PAGE DIMENSIONS
\usepackage[margin=25mm,inner=30mm]{geometry} % to change the page dimensions


%%% HEADERS & FOOTERS   %%%
\usepackage{fancyhdr} 	% This should be set AFTER setting up the page geometry
\pagestyle{fancy} 			% options: empty , plain , fancy
\renewcommand{\headrulewidth}{0.5pt} % customise the layout...

\renewcommand{\footrulewidth}{0.5pt} % customise the layout...
\lhead{\includegraphics[scale=.15]{Logos/header_utfsm.png}}\cfoot{}
\cfoot{\textbf{\thepage}}
\rfoot{\therevision}
\lfoot{\theAssignment}\chead{}
\rhead{}


\usepackage[colorlinks=true,linkcolor=blue,urlcolor=blue]{hyperref}

%%% SECTION TITLE APPEARANCE
\usepackage{sectsty}
%\allsectionsfont{\sffamily\mdseries\upshape} % (See the fntguide.pdf for font help)
\usepackage{float}
\usepackage{titlesec}
\titlelabel{\thetitle.\quad}

\usepackage[numbers]{natbib}

\usepackage{lipsum}
\usepackage[section]{placeins}

%%% ESTILO DE CÓDIGO SQL
\lstdefinestyle{sqlstyle}{
  language=SQL,
  basicstyle=\ttfamily\small,
  keywordstyle=\bfseries\color{blue!60!black},
  stringstyle=\color{orange!70!black},
  commentstyle=\color{gray},
  numbers=none,
  columns=flexible,
  keepspaces=true,
  upquote=true,
  showstringspaces=false,
  breaklines=true,
  frame=single,
  rulecolor=\color{gray!60},
  backgroundcolor=\color{gray!8},
  framesep=4pt,
  xleftmargin=2pt,
  xrightmargin=2pt,
  aboveskip=6pt,
  belowskip=6pt,
  % tildes y ñ dentro del código (p. ej. 'Montaña'): listings no lee UTF-8 por sí solo
  literate={á}{{\'a}}1 {é}{{\'e}}1 {í}{{\'i}}1 {ó}{{\'o}}1 {ú}{{\'u}}1
           {Á}{{\'A}}1 {É}{{\'E}}1 {Í}{{\'I}}1 {Ó}{{\'O}}1 {Ú}{{\'U}}1
           {ñ}{{\~n}}1 {Ñ}{{\~N}}1 {ü}{{\"u}}1
}
\lstset{style=sqlstyle}

% Figura de captura de MySQL Workbench (ancho común para todas)
\newcommand{\captura}[3]{%
  \begin{figure}[H]
    \centering
    \includegraphics[width=0.88\linewidth]{Lab5/Figuras/#1}
    \caption{#2}
    \label{#3}
  \end{figure}
}


%------------------------------------------------------------------------------
%%%% TITLE PAGE
%------------------------------------------------------------------------------
\begin{document}
\selectlanguage{spanish}

\begin{titlepage}
\begin{figure}[H]
\includegraphics[width = 0.4\linewidth]{Logos/logo_di.png}
\hspace{144pt}
\vspace{-70pt}
\end{figure}
\vspace{60mm}
\begin{center}
\hrule
\vspace{2mm}
{\Huge \textbf{\sc Informe de Laboratorio 5:}} \\
\vspace{2mm}
{\Large \sc{\textit{Entregable SQL (2/2): base de datos de Andes Bike}}}
\vspace{2mm}
\hrule
\vspace{1cm}


%%%%AUTHORS%%%%
\centering
    {\textbf{Diego Victoriano Rodríguez} \\
    {\small Ingeniería Civil Industrial, 21.540.820-5, \emph{diego.victoriano@usm.cl}}\\
    }

    \vspace{0.5cm}
    {ICN292 - Sistemas de Información para la Gestión.} \\
    {\textbf{Paralelo:} 101}

    \vspace{0.5cm}
    \textbf{Profesores:} José Miguel González P. \textbar\ José Luis Sáez T. \textbar\ Patricio Cifuentes\\[0.3cm]
    \textbf{Ayudantes:} Valeria Bobadilla \textbar  Catalina Bodaleo \textbar\ Luis Contreras \textbar  Benjamín Jara

    \vspace{0.8cm}
    \textbf{Repositorio GitHub:}\\[0.1cm]
    {\small \url{https://github.com/diegovictorianor/ICN292-Lab5-Victoriano-Diego}}

    \vfill
    \textsc{\monthyeardate\today}\\[0.5cm]
    \textsc{\theCampus}
\end{center}


\clearpage

\parindent=0mm


\end{titlepage}

%--------------------------------------------------
%%% CONSULTAS %%%%
%--------------------------------------------------
%%% Consulta 1
\section*{Consulta 1: vendedores con su cantidad de ventas}

\begin{lstlisting}
SELECT v.nombre,
       COUNT(ve.id_venta) AS cantidad_ventas
FROM vendedores v
LEFT JOIN ventas ve ON v.id_vendedor = ve.id_vendedor
GROUP BY v.id_vendedor, v.nombre
ORDER BY cantidad_ventas DESC, v.nombre;
\end{lstlisting}

\captura{consulta_1.png}{Resultado de la consulta 1 en MySQL Workbench.}{fig:c1}

\clearpage
%%% Consulta 2
\section*{Consulta 2: ventas y unidades vendidas por tienda}
 
\begin{lstlisting}
SELECT t.nombre,
       COUNT(ve.id_venta)             AS total_ventas,
       IFNULL(SUM(ve.unidades), 0)    AS total_unidades
FROM tiendas t
LEFT JOIN vendedores v  ON t.id_tienda   = v.id_tienda
LEFT JOIN ventas     ve ON v.id_vendedor = ve.id_vendedor
GROUP BY t.id_tienda, t.nombre
ORDER BY total_unidades DESC, t.nombre;
\end{lstlisting}
 
\captura{consulta_2.png}{Resultado de la consulta 2 en MySQL Workbench.}{fig:c2}

\clearpage
%%% Consulta 3
\section*{Consulta 3: tiendas sin ningún vendedor asignado}

\begin{lstlisting}
SELECT t.nombre, t.comuna
FROM tiendas t
LEFT JOIN vendedores v ON t.id_tienda = v.id_tienda
WHERE v.id_vendedor IS NULL;
\end{lstlisting}

\captura{consulta_3.png}{Resultado de la consulta 3 en MySQL Workbench.}{fig:c3}

\clearpage
%%% Consulta 4
\section*{Consulta 4: modelos sin ventas en junio de 2025}

\begin{lstlisting}
SELECT b.modelo, b.categoria
FROM bicicletas b
LEFT JOIN ventas ve ON b.id_bicicleta = ve.id_bicicleta
                   AND ve.fecha >= '2025-06-01'
                   AND ve.fecha <  '2025-07-01'
WHERE ve.id_venta IS NULL
ORDER BY b.modelo;
\end{lstlisting}

La condición del mes viene en el \texttt{ON} del \texttt{LEFT JOIN}. Si la hubiera puesto en el \texttt{WHERE}, ese filtro descartaría las filas sin pareja (con \texttt{fecha} NULL) y la consulta devolvería cero filas.

\captura{consulta_4.png}{Resultado de la consulta 4 en MySQL Workbench.}{fig:c4}

\clearpage
%%% Consulta 5
\section*{Consulta 5: vendedores sin ventas y su tienda}

\begin{lstlisting}
SELECT v.nombre AS vendedor,
       t.nombre AS tienda
FROM vendedores v
LEFT JOIN tiendas t  ON v.id_tienda   = t.id_tienda
LEFT JOIN ventas  ve ON v.id_vendedor = ve.id_vendedor
WHERE ve.id_venta IS NULL
ORDER BY v.nombre;
\end{lstlisting}

\captura{consulta_5.png}{Resultado de la consulta 5 en MySQL Workbench.}{fig:c5}

\clearpage
%%% Consulta 6
\section*{Consulta 6: pares de vendedores de la misma tienda}

\begin{lstlisting}
SELECT v1.nombre        AS vendedor_1,
       v1.fecha_ingreso AS ingreso_1,
       v2.nombre        AS vendedor_2,
       v2.fecha_ingreso AS ingreso_2,
       t.nombre         AS tienda
FROM vendedores v1
JOIN vendedores v2 ON v1.id_tienda = v2.id_tienda
                  AND (v1.fecha_ingreso < v2.fecha_ingreso
                       OR (v1.fecha_ingreso = v2.fecha_ingreso
                           AND v1.id_vendedor < v2.id_vendedor))
JOIN tiendas t ON v1.id_tienda = t.id_tienda
ORDER BY t.nombre, v1.fecha_ingreso, v2.fecha_ingreso;
\end{lstlisting}

\captura{consulta_6.png}{Resultado de la consulta 6 en MySQL Workbench.}{fig:c6}

\clearpage
%%% Consulta 7
\section*{Consulta 7: monto vendido en Montaña y Ruta por tienda}

\begin{lstlisting}
SELECT t.nombre,
       SUM(ve.unidades * b.precio) AS monto_total
FROM tiendas t
JOIN vendedores v  ON t.id_tienda    = v.id_tienda
JOIN ventas     ve ON v.id_vendedor  = ve.id_vendedor
JOIN bicicletas b  ON ve.id_bicicleta = b.id_bicicleta
WHERE b.categoria IN ('Montaña', 'Ruta')
GROUP BY t.id_tienda, t.nombre
HAVING SUM(ve.unidades * b.precio) > 2000000
ORDER BY monto_total DESC;
\end{lstlisting}

\captura{consulta_7.png}{Resultado de la consulta 7 en MySQL Workbench.}{fig:c7}

\clearpage
%%% Consulta 8
\section*{Consulta 8: bicicletas sobre el precio promedio del catálogo}

\begin{lstlisting}
SELECT modelo, precio
FROM bicicletas
WHERE precio > (SELECT AVG(precio) FROM bicicletas)
ORDER BY precio DESC;
\end{lstlisting}

El promedio del catálogo es \$657.490 (verificado en la Figura~\ref{fig:c8b}). en esta consulta se obtuvo mediante la subconsulta \lstinline|SELECT AVG(precio) FROM bicicletas|, que promedia el precio de los 8 modelos de la tabla \texttt{bicicletas}.

\captura{consulta_8.png}{Resultado de la consulta 8 en MySQL Workbench.}{fig:c8}

\clearpage
\subsection*{Verificación del promedio del catálogo}

\begin{lstlisting}
SELECT AVG(precio) AS promedio_catalogo
FROM bicicletas;
\end{lstlisting}

\captura{verificacion_promedio.png}{Verificación del promedio del catálogo en MySQL Workbench.}{fig:c8b}

\clearpage
%%% Consulta 9
\section*{Consulta 9: vendedores sobre el promedio de unidades de la cadena}

\begin{lstlisting}
SELECT v.nombre,
       AVG(ve.unidades) AS promedio_unidades
FROM vendedores v
JOIN ventas ve ON v.id_vendedor = ve.id_vendedor
GROUP BY v.id_vendedor, v.nombre
HAVING AVG(ve.unidades) > (SELECT AVG(unidades) FROM ventas)
ORDER BY promedio_unidades DESC, v.nombre;
\end{lstlisting}

\captura{consulta_9.png}{Resultado de la consulta 9 en MySQL Workbench.}{fig:c9}

\clearpage
%%% Consulta 10
\section*{Consulta 10: monto total vendido por vendedor}

\begin{lstlisting}
SELECT v.nombre,
       SUM(ve.unidades * b.precio) AS monto_total
FROM vendedores v
LEFT JOIN ventas     ve ON v.id_vendedor    = ve.id_vendedor
LEFT JOIN bicicletas b  ON ve.id_bicicleta  = b.id_bicicleta
GROUP BY v.id_vendedor, v.nombre
ORDER BY monto_total DESC, v.nombre;
\end{lstlisting}

Fernanda Lillo y Rubén Castro muestran NULL y no cero porque, al no tener ventas, el \texttt{LEFT JOIN} les deja \texttt{unidades} y \texttt{precio} en NULL, y \texttt{SUM} de puros NULL devuelve NULL. Para que mostraran cero habría que envolver la suma en \mbox{\lstinline|IFNULL(SUM(ve.unidades * b.precio), 0)|}.

\captura{consulta_10.png}{Resultado de la consulta 10 en MySQL Workbench.}{fig:c10}

\end{document}
